\section{Introduction}
\label{sec:introduction}

Integer division and modulo operations are central tools in discrete
mathematics, number theory, and algorithms. While their properties are well
known, rigorous formalization and verification, particularly via recursive
definitions, offer an interesting alternative to the traditional axiomatic
model.

This article takes the recursive route. Instead of assuming native division
and modulo, it defines a state $DivMod(a,b,q,r)$ where $a = bq + r$, then proves
that normalizing the pair $(q,r)$ preserves the represented dividend and
reaches the canonical remainder interval. The familiar operations
$\operatorname{div}$ and $\operatorname{mod}$ are then projections from that
normalized state.

The mathematical statements below are backed by Scala source verified with
\href{https://epfl-lara.github.io/stainless/intro.html}{Stainless}. The article
keeps the proof discussion centered on the properties; source links point to
the maintained verification code.

This article establishes:

\begin{itemize}
  \item Foundational identities: trivial case, self-identity, division by one,
    and agreement with the native modulo operator (Sections~6.1--6.4).
  \item Linear shift laws under single-step and multiple-step divisor addition
    (Sections~6.5--6.6).
  \item Uniqueness and idempotence of the normalized remainder
    (Sections~6.7--6.8).
  \item Distributivity of modulo and division over addition and subtraction
    (Sections~6.9--6.10).
  \item Divisible-base shift invariance and symmetric remainder pairs
    (Sections~6.11--6.12).
  \item The unit-step increment law and zero-density over consecutive integers
    (Sections~6.13--6.14).
\end{itemize}

\section{Limitations}
\label{sec:limitations}

The implementation presented in this article is limited to the division and
modulo operations for integers. Its goal is to make available a set of lemmas
and proofs that can be verified and used as a base to prove other properties
related to the division and modulo operations. Therefore, the implementation
is optimized to correctness and not to performance.

The use of \texttt{BigInt} in the implementation focused on unbounded
integers, without the need to worry about overflow or underflow issues. But
they are still constrained by the memory available in the system. Similarly,
some lemmas and proofs use the recursive definition of the division and modulo
operations, which could trigger a stack overflow for large numbers. Those
issues do not invalidate the mathematical properties proved in this article,
which are the main focus of this article.
